\chapter{연결이란 무엇인가}\label{ch:connect}

이 문서에는 기계가 셋 나온다. 드럼이 도는 기어식, 릴레이로 된 릴레이식, 움직이는 부품이 없는 FPGA다. 생김새는
전혀 다르지만 세 기계가 하는 일은 하나다. 앞의 두 장(\ref{ch:enigma-math}장, \ref{ch:cryptanalysis}장)은 그 일을
식으로 보였다. 이 장과 다음 장은 같은 일을 코드도 회로도 없이 \emph{전선}으로 다시 설명한다. 이 장은 앞 장의 추론이
왜 ``선과 선의 연결''이라는 모양이 되는지를 본다. 다음 장(\ref{ch:find}장)은 그 연결을 따라가는 방법을 보고, 세
기계가 그 방법을 각각 어떻게 하는지로 이 문서 전체의 지도를 만든다.

\section{에니그마, 크립, 메뉴}\label{sec:bombe-one-page}

앞 두 장에서 이 장이 쓰는 것만 다시 적는다.

\begin{itemize}
\item 크립 위치 $k$의 평문 글자 $p$는 플러그보드 $\Stk$, \term{스크램블러}{scrambler}(로터와 반사판) $\Scr_k$, 다시 $\Stk$를
지나 $c = \Stk(\Scr_k(\Stk(p)))$가 된다(식~\eqref{eq:am-enigma}). 스크램블러는 되돌릴 수 있다. $x$를 $y$로 바꾸는 순간이면
$y$를 넣어도 $x$가 나온다. 플러그보드는 글자 쌍을 맞바꾸는 배선판이고, 서로 이어진 두 글자를 \term{스테커}{Stecker}
짝이라 한다.
\item 플러그보드를 한쪽으로 넘기면 $\Stk(c) = \Scr_k(\Stk(p))$다(식~\eqref{eq:ac-moved}). ``$p$의 짝이 $x$라면 $c$의
짝은 $\Scr_k(x)$다''라는 추론이고, 거꾸로도 읽힌다.
\item 크립의 글자 관계를 선으로 그린 그래프가 \term{메뉴}{menu}다(그림~\ref{fig:g-menu}). 걸러 내는 힘은 출발한 글자로
돌아오는 \term{루프}{loop}와, 플러그보드의 대칭에서 나온다(\ref{sec:ac-loop}절, \ref{sec:ac-diagonal}절).
\end{itemize}

봄베는 그날의 로터 위치를 모르므로 로터 순서 하나에 대해 위치 $26^3 = 17{,}576$가지를 하나씩 지나간다. 위치마다
묻는 것은 하나다. \emph{``이 위치에서 메뉴의 모든 관계를 동시에 만족하는 플러그보드가 있을 수 있는가?''} 있을 수
없으면 다음 위치로 가고, 있을 수 있으면 멈춘다. 멈추는 것이 \term{정지}{stop}다(\ref{sec:ac-stop}절).

그림~\ref{fig:g-bombe-front}는 1940년대의 영국 봄베를 문헌의 기술에 따라 그린 것이다. 앞면의 드럼 세 개
한 벌이 스크램블러 하나이고, 뒷면의 잭에 케이블을 꽂아 메뉴를 배선한다. 이 장의 나머지는 이 기계가 위의 물음을
어떻게 ``전선의 연결''로 바꾸었는지를 본다.

\begin{figure}[p]
\centering
% 영국 3로터 튜링-웰치먼 봄베의 앞면(위)과 뒷면 잭 배치(아래). 직접 그린 도해이며 사진을 옮긴 것이 아니다.
% 근거: docs/reference/bombe-mechanics.md 2, 5, 11절 (us6812report1944, ellsbury1998bombe, jackson-bh10).
% 단위: 1 = 1 ft. 근거가 있는 것: 뱅크 3개 x 유닛 12개 x 드럼 3개(108개), 지시 드럼 3개가 가운데 뱅크
% 오른쪽, 외형 약 7 ft x 6 ft 6 in, 드럼 지름 약 5 in, 뒷면 잭 열 9개의 용도.
% 어림해 그린 것: 드럼 사이 간격과 캐비닛 안에서의 정확한 위치, 지시 드럼의 배열 방향, 바퀴, 잭의 개수와 간격.
\begin{tikzpicture}[x=1.27cm,y=1.27cm,
    drum/.style={circle,draw=black!70,fill=clight,minimum size=5.3mm,inner sep=0pt,line width=0.4pt},
    lab/.style={font=\footnotesize,align=left,anchor=west},
    lead/.style={draw=cgray,line width=0.4pt}]
  % ---------------- 앞면 ----------------
  \draw[line width=0.9pt,fill=white] (0,0) rectangle (7,6.5);
  \foreach \x in {0.5,6.5} { \draw[fill=cgray!40] (\x,-0.12) circle (0.1); }
  % 드럼: 뱅크 3개 x 3줄 x 12열
  \foreach \b/\by in {0/4.55,1/2.8,2/1.05} {
    \fill[cbluebg!50] (0.33,\by-0.25) rectangle (6.02,\by+1.25);
    \foreach \r in {0,1,2} {
      \foreach \c in {0,...,11} {
        \node[drum] at (0.6+\c*0.47,\by+1.0-\r*0.5) {};
        \fill[black!70] (0.6+\c*0.47,\by+1.0-\r*0.5) circle (0.03);
      }
    }
  }
  % 스크램블러 한 벌 강조 (위 뱅크, 셋째 열)
  \draw[corange,line width=1.1pt,rounded corners=2pt] (1.32,4.32) rectangle (1.76,5.78);
  % 지시 드럼 (가운데 뱅크 오른쪽)
  \foreach \r in {0,1,2} { \node[drum,fill=cyellowbg] (ind\r) at (6.5,3.8-\r*0.5) {}; }
  % 치수
  \draw[<->,cgray] (0,6.8) -- node[above,font=\footnotesize,text=cgray] {폭 약 7 ft (2.1 m)} (7,6.8);
  \draw[<->,cgray] (-0.35,0) -- node[above,rotate=90,font=\footnotesize,text=cgray] {높이 약 6 ft 6 in (2.0 m)} (-0.35,6.5);
  % 라벨
  \draw[lead] (6.02,5.3) -- (7.3,5.9);
  \node[lab] at (7.3,5.9) {\textbf{뱅크}(bank) × 3\\유닛 12개씩, 모두 36개};
  \draw[lead,corange] (1.54,5.78) -- (1.9,6.15) -- (6.5,6.15) -- (7.3,4.95);
  \node[lab,text=corange] at (7.3,4.75) {\textbf{스크램블러 한 벌}(유닛)\\세로로 놓인 드럼 3개\\= 플러그보드 없는 에니그마};
  \draw[lead] (ind0.east) -- (7.3,3.5);
  \node[lab] at (7.3,3.5) {\textbf{지시 드럼}(indicator drums) 3개\\정지한 위치를 읽는다};
  \node[lab] at (7.3,2.1) {한 벌 안의 드럼 (위에서 아래로)\\위: 연속 회전 --- 왼쪽 로터\\가운데 --- 가운데 로터\\아래: 가장 느림 --- 오른쪽 로터};
  \node[lab] at (7.3,0.75) {드럼 108개, 지름 약 5 in\\깊이 약 2 ft, 무게 약 1톤};
  \node[font=\small\bfseries,anchor=west] at (0,7.35) {앞면};
  % ---------------- 뒷면 (잭 배치 도식) ----------------
  \begin{scope}[yshift=-6.6cm]
  \node[font=\small\bfseries,anchor=west] at (0,3.75) {뒷면: 메뉴를 배선하는 잭 (배치 도식)};
  \draw[line width=0.9pt,fill=white] (0,0) rectangle (7,3.4);
  \foreach \g in {0,1,2} {
    \draw[cgray,dashed] (0.25+\g*2.25,0.15) rectangle (2.25+\g*2.25,3.25);
    % 대각선 보드 열
    \foreach \j in {0,...,12} { \fill[cblue] (0.55+\g*2.25,0.4+\j*0.215) rectangle ++(0.22,0.12); }
    % commons 열
    \foreach \j in {0,...,7} { \fill[cgreen] (1.14+\g*2.25,0.5+\j*0.34) rectangle ++(0.22,0.2); }
    % 스크램블러 IN/OUT 열 (유닛당 한 쌍, 12쌍)
    \foreach \j in {0,...,11} {
      \fill[corange] (1.68+\g*2.25,0.4+\j*0.235) rectangle ++(0.13,0.13);
      \fill[corange!55] (1.84+\g*2.25,0.4+\j*0.235) rectangle ++(0.13,0.13);
    }
  }
  % 메뉴 케이블 예 (개념)
  \draw[black!75,line width=0.9pt] (1.75,2.8) .. controls (1.2,2.3) and (1.0,1.9) .. (0.66,1.55);
  \draw[black!75,line width=0.9pt] (1.91,2.8) .. controls (2.3,2.0) and (1.6,1.4) .. (1.25,1.28);
  \draw[black!75,line width=0.9pt] (1.25,1.6) .. controls (0.9,1.2) and (0.9,0.9) .. (0.66,0.68);
  \draw[lead,cblue] (0.66,3.0) -- (0.66,3.5) -- (7.3,3.1);
  \node[lab,text=cblue] at (7.3,3.1) {\textbf{대각선 보드}(diagonal board)\\글자마다 26접점 잭, 뱅크마다 한 벌};
  \draw[lead,cgreen] (5.75,2.0) -- (7.3,2.0);
  \node[lab,text=cgreen] at (7.3,2.0) {\textbf{커먼}(commons)\\한 글자에 선이 여럿 모일 때의 병렬 잭};
  \draw[lead,corange] (6.47,1.0) -- (7.3,0.9);
  \node[lab,text=corange] at (7.3,0.9) {\textbf{스크램블러 IN/OUT 잭}\\유닛마다 한 쌍. 26선 케이블로 잇는다};
  \end{scope}
\end{tikzpicture}

\caption{영국 3로터 봄베의 도해(직접 그린 것이며 사진을 옮긴 것이 아니다). 뱅크와 드럼의 수, 지시 드럼의
위치, 뒷면 잭 열의 용도, 외형 치수는 문헌~\cite{us6812report1944,ellsbury1998bombe,jackson-bh10}을
따랐다. 드럼의 정확한 간격과 위치, 지시 드럼의 배열 방향, 잭의 개수와 간격은 어림해 그렸다.}
\label{fig:g-bombe-front}
\end{figure}

\section{가설 하나를 선 하나로}\label{sec:c-wire}

\Why 플러그보드를 모른다는 것은 글자마다 짝을 모른다는 것이다. 글자 \texttt{E}의 짝은 \texttt{A}일 수도,
\texttt{B}일 수도, \ldots, \texttt{Z}일 수도 있다(짝이 없으면 자기 자신이다). 가능성이 26가지이니 \texttt{E}에
대한 가설은 26개다. 봄베는 이 가설 하나하나에 전선 한 가닥을 준다. 글자마다 26가닥을 묶은
\term{케이블}{cable}이 있고, 케이블 \texttt{E}의 \texttt{v}번 선에 전류가 흐른다는 것은 ``\texttt{E}의 짝이
\texttt{V}다''라는 가설이 \emph{살아 있다}는 뜻이다(그림~\ref{fig:c-cable}). 글자가 26개이므로 선은 모두
$26 \times 26 = 676$가닥이다.

\begin{figure}[ht]
\centering
\begin{tikzpicture}[x=4.6mm,y=1cm]
  \node[font=\small,anchor=east] at (-1.0,0) {케이블 \texttt{E}};
  \foreach \l [count=\i from 0] in {a,b,c,d,e,f,g,h,i,j,k,l,m,n,o,p,q,r,s,t,u,v,w,x,y,z} {
    \draw[cgray!70,fill=white] (\i-0.42,-0.3) rectangle (\i+0.42,0.3);
    \node[font=\ttfamily\scriptsize] at (\i,-0.55) {\l};
  }
  \fill[cblue] (21-0.42,-0.3) rectangle (21+0.42,0.3);
  \draw[<-,cblue] (21,0.35) -- (21,0.9) node[above,font=\footnotesize] {``\texttt{E}의 짝은 \texttt{V}다''에 전류};
  \draw[decorate,decoration={brace,amplitude=4pt,mirror}] (-0.42,-0.8) -- (25.42,-0.8)
    node[midway,below=3pt,note] {가설 26개 = 선 26가닥. 짝은 하나뿐이므로 이 가운데 참은 꼭 하나다.};
\end{tikzpicture}
\caption{케이블 하나는 글자 하나에 대한 가설 26개다. 케이블이 26개이므로 봄베의 선은 676가닥이다.}
\label{fig:c-cable}
\end{figure}

\General 퀴즈의 보기에 번호를 매기고 보기마다 전구를 하나씩 단다고 생각하면 된다. 전구가 켜진 보기는 ``아직
지워지지 않은 보기''다. 정답은 하나뿐이므로, 어느 문제의 전구가 \emph{모두} 켜졌다면 무엇인가 앞에서 잘못된
것이다.

\Ours 이 676가닥이 세 기계의 상태 전부다. 기어식과 릴레이식에서는 실제 전선이고, FPGA와 Rust에서는 676개의
비트다(\ref{ch:rust-bombe}장의 \code{live} 배열, \ref{ch:rtl-core}장의 \code{live\_q} 레지스터). 무엇으로
만들든 ``선이 켜졌다 = 가설이 살아 있다''는 뜻은 같다.

\section{전기에서 빌려 올 사실 하나}\label{sec:c-current}

이 두 장에서 필요한 전기 지식은 한 줄이다. \emph{전류는 끊기지 않은 길로만 흐르고, 전압을 건 선과 이어진
선은 모두 같은 전압이 된다.} 두 선 사이에 금속이 닿아 있으면 이어져 있고, 스위치가 열려 있으면 끊겨 있다.
이어진 선이 다시 다른 선과 이어져 있으면 그 선도 같은 전압이 된다. 전압·전류·저항은 \ref{ch:39-electric-basics}장에서
제대로 다룬다.

이 문서에서 \term{연결}{connection}은 이 뜻이다. 두 선이 연결되어 있으면 한쪽에 전압을 걸 때 다른 쪽에도
전압이 온다. 거꾸로도 마찬가지다. 전선의 연결에는 방향이 없다.

\section{스크램블러: 추론 하나가 연결 하나}\label{sec:c-scrambler}

\Why 메뉴의 선 하나, 곧 ``$k$번째 글자에서 평문 $p$가 암호문 $c$가 되었다''를 생각한다. 에니그마 안에서 전류는
플러그보드를 지나 $p$의 짝이 되고, 스크램블러를 지나고, 다시 플러그보드를 지나 $c$가 되었다. 플러그보드는
되돌릴 수 있으므로 이것은 ``$p$의 짝을 스크램블러에 넣으면 $c$의 짝이 나온다''와 같다. 그러면 추론이 하나
생긴다.

\begin{center}
``$p$의 짝이 $x$라면, $c$의 짝은 스크램블러가 $x$를 바꾼 글자다.''
\end{center}

봄베는 이 추론을 전선으로 만든다. 메뉴의 선마다 스크램블러 하나를 두고, 케이블 $p$의 $x$번 선을 그
스크램블러를 지나 케이블 $c$의 (바뀐 글자)번 선에 \emph{잇는다}. 26가닥 모두에 대해서다. 스크램블러가
되돌릴 수 있으므로 이 연결은 거꾸로 읽어도 맞는 추론이다. ``$c$의 짝이 $y$라면 $p$의 짝은 스크램블러가 $y$를
바꾼 글자다.'' 그래서 방향이 없는 전선으로 만들어도 된다.

그림~\ref{fig:c-scr}는 글자가 여섯 개(\texttt{A}--\texttt{F})뿐인 작은 봄베에서 스크램블러 하나가 두 케이블을
잇는 모습이다. 이 장과 다음 장의 그림은 모두 이 작은 봄베로 그린다. 글자가 26개이면 선이 많아질 뿐 원리는 같다.

\begin{figure}[ht]
\centering
\begin{tikzpicture}[x=1cm,y=-5.5mm,
    dot/.style={circle,draw,fill=white,inner sep=1.5pt}]
  \foreach \col/\name in {0/A,3.6/B} {
    \node[font=\small\bfseries] at (\col,-1.1) {케이블 \texttt{\name}};
    \foreach \l [count=\i from 0] in {a,b,c,d,e,f} {
      \node[dot] (\name\l) at (\col,\i) {};
    }
  }
  \foreach \l [count=\i from 0] in {a,b,c,d,e,f} {
    \node[font=\ttfamily\scriptsize,anchor=east] at (-0.15,\i) {\l};
    \node[font=\ttfamily\scriptsize,anchor=west] at (3.75,\i) {\l};
  }
  \foreach \x/\y in {a/b,b/a,c/f,d/e,e/d,f/c} { \draw[wire] (A\x) -- (B\y); }
  \draw[live] (Ad) -- (Be);
  \node[blkn,font=\footnotesize] at (1.8,6.3) {스크램블러 $\Scr_1$\\{\scriptsize a$\leftrightarrow$b, c$\leftrightarrow$f, d$\leftrightarrow$e}};
  \node[note,anchor=west,align=left] at (5.0,1.5) {파란 선: ``\texttt{A}의 짝이 \texttt{D}이면\\
    \texttt{B}의 짝은 \texttt{E}다''이자\\``\texttt{B}의 짝이 \texttt{E}이면\\\texttt{A}의 짝은 \texttt{D}다''};
\end{tikzpicture}
\caption{작은 봄베(글자 여섯 개)에서 메뉴의 선 \texttt{A}--\texttt{B} 하나. 스크램블러가 바꾸는 쌍대로 두 케이블의
선이 이어진다. 선 하나가 추론 하나이고, 방향이 없으므로 양쪽으로 읽힌다.}
\label{fig:c-scr}
\end{figure}

\General 지하철 노선도의 환승 통로와 비슷하다. 통로는 두 역을 잇고, 어느 쪽에서 들어가도 다른 쪽으로 나온다.
노선도를 보고 ``이 역에서 갈 수 있는 역''을 묻는 것과 봄베의 물음은 같은 꼴이다.

\Ours 실제 봄베에서는 드럼 세 개 한 벌이 스크램블러 하나이고, 메뉴의 선 하나마다 한 벌을 쓴다. 크립이 16글자이고
그 모든 관계를 메뉴에 쓰면 메뉴의 선이 16개이고, 스크램블러 16벌이 $16 \times 26 = 416$개의 연결을 만든다.

\section{대각선 보드: 늘 있는 연결}\label{sec:c-diagonal}

\Why 플러그보드는 글자를 \emph{쌍으로} 바꾼다. 그러므로 ``\texttt{A}의 짝이 \texttt{C}다''와 ``\texttt{C}의 짝이
\texttt{A}다''는 같은 말이다. 하나가 살아 있으면 다른 하나도 살아 있어야 한다. 이것을 전선으로 넣은 것이
\term{대각선 보드}{diagonal board}다. 케이블 $X$의 $y$번 선과 케이블 $Y$의 $x$번 선을 잇는다.

이름은 그림~\ref{fig:c-diag}에서 나왔다. 676가닥을 케이블(가로줄) × 선(세로줄)의 정사각형으로 늘어놓으면,
대각선 보드가 잇는 두 칸은 대각선을 사이에 두고 마주 보는 거울상 자리다. 대각선 위의 칸(``\texttt{A}의 짝은
\texttt{A}다'', 곧 짝 없음)은 자기 자신과 짝이 된다. 26글자이면 대각선 밖의 칸 $676 - 26 = 650$개가 둘씩 묶여
연결 325개가 된다.

\begin{figure}[ht]
\centering
\begin{tikzpicture}[x=7mm,y=-7mm]
  \foreach \r/\R in {0/A,1/B,2/C,3/D,4/E,5/F} {
    \node[font=\ttfamily\footnotesize,anchor=east] at (-0.6,\r) {\R};
    \foreach \c in {0,...,5} {
      \ifnum\r=\c \def\fillc{clight}\else \def\fillc{white}\fi
      \draw[cgray!70,fill=\fillc] (\c-0.45,\r-0.45) rectangle (\c+0.45,\r+0.45);
    }
  }
  \foreach \l [count=\c from 0] in {a,b,c,d,e,f} { \node[font=\ttfamily\footnotesize] at (\c,-0.9) {\l}; }
  \node[font=\footnotesize,anchor=east] at (-1.2,2.5) {케이블};
  \node[font=\footnotesize] at (2.5,-1.6) {선};
  \draw[cgray,dashed] (-0.45,-0.45) -- (5.45,5.45);
  \fill[cblue] (2-0.45,0-0.45) rectangle (2+0.45,0+0.45);
  \fill[cblue] (0-0.45,2-0.45) rectangle (0+0.45,2+0.45);
  \draw[live,<->,bend left=35] (2.2,0.35) to (0.35,2.2);
  \fill[corange] (4-0.45,1-0.45) rectangle (4+0.45,1+0.45);
  \fill[corange] (1-0.45,4-0.45) rectangle (1+0.45,4+0.45);
  \draw[line width=1.1pt,corange,<->,bend left=25] (4.1,1.4) to (1.4,4.1);
  \node[note,anchor=west,align=left] at (6.4,0.6) {파란 칸: ``\texttt{A}의 짝은 \texttt{C}''와\\``\texttt{C}의 짝은 \texttt{A}''};
  \node[note,anchor=west,align=left] at (6.4,2.4) {주황 칸: ``\texttt{B}의 짝은 \texttt{E}''와\\``\texttt{E}의 짝은 \texttt{B}''};
  \node[note,anchor=west,align=left] at (6.4,4.2) {회색 칸(대각선): 짝 없음.\\자기 자신과 이어진다};
\end{tikzpicture}
\caption{대각선 보드. 선 전체를 케이블 × 선의 정사각형으로 놓으면, 이어지는 두 선은 대각선을 사이에 둔 거울상
자리다. 작은 봄베에서는 연결이 15개, 26글자에서는 325개다.}
\label{fig:c-diag}
\end{figure}

\General 친구 관계와 같다. ``철수가 영희의 친구다''가 참이면 ``영희가 철수의 친구다''도 참이다. 친구 명단을
표로 만들면 표가 대각선에 대해 대칭이 된다.

\Ours 대각선 보드의 연결은 드럼 위치와 무관하다. 기계가 도는 동안 한 번도 바뀌지 않는 고정 배선이다. 실제
봄베에서도 그랬고(웰치먼의 개량), 이 프로젝트의 세 기계에서도 그렇다. 릴레이식과 기어식은 함께 쓰는 기판 다섯
장의 구리 배선으로, FPGA는 스위치도 선택도 없는 고정 배선으로(\ref{ch:fpga-engine}장) 만든다. 메뉴의 선만으로도 봄베는 동작하지만, 이 연결이 더해지면 모순이 훨씬 빨리 드러나
쓸모없는 정지가 크게 줄어든다.

\section{위치가 바뀌면 연결이 바뀐다}\label{sec:c-position}

스크램블러의 치환은 로터의 위치에 따라 달라진다. 그러므로 메뉴의 선이 만드는 연결은 드럼 위치마다 다르다.
게다가 메뉴의 선마다 크립 안의 위치 $k$가 다르므로, 같은 순간에도 선마다 $k$칸씩 어긋난 위치의 스크램블러를 쓴다.
메뉴의 선이 16개라면, 봄베가 위치를 하나 넘길 때마다 416개의 연결이 한꺼번에 다른 연결로 바뀐다. 대각선 보드의
325개만 그대로다.

그래서 봄베의 일은 두 가지로 나뉜다.

\begin{keybox}[title={봄베의 두 가지 일}]
\begin{enumerate}
\item \textbf{연결을 만든다.} 지금 위치에서 메뉴의 선마다 스크램블러가 어떤 연결을 만드는지 정한다. 위치를 넘기면
다시 만든다.
\item \textbf{이어진 선을 찾는다.} 가설 하나의 선에 전압을 걸고, 그 선과 (여러 연결을 거쳐) 이어진 선을 모두 찾는다.
\end{enumerate}
\end{keybox}

세 기계는 이 두 일을 서로 다른 수단으로 한다. 첫째 일만 미리 보면 이렇다.

\begin{itemize}
\item \textbf{기어식}은 실제 봄베처럼 드럼을 돌린다. 드럼 안에 로터의 배선이 들어 있어, 드럼이 한 칸 돌면 배선이
닿는 자리가 바뀌고 연결이 바뀐다(\ref{ch:50-gear-mechanics}장).
\item \textbf{릴레이식}에는 도는 것이 없다. 로터의 회전을 ``몇 칸 밀어서 배선에 넣는가''로 바꾸고, 그 밀기를
릴레이 접점이 고른다(\ref{ch:relay-machine}장).
\item \textbf{FPGA}에는 움직이는 것도 접점도 없다. 치환표를 메모리에 두고 위치에 맞는 표를 골라 읽는다
(\ref{ch:rtl-scrambler}장, \ref{ch:fpga-engine}장). Rust의 기준 모델도 같다(\ref{ch:rust-bombe}장).
\end{itemize}

\section{연결은 그래프다}\label{sec:c-graph}

선을 점으로, 연결을 점 사이의 선분으로 그리면 \term{그래프}{graph}가 된다. 메뉴도 그래프였지만(글자가 점) 이것은
한 층 아래의 그래프다. 점이 676개이고, 선분은 메뉴가 만드는 연결과 대각선 보드의 연결이다.

그림~\ref{fig:c-graph}는 작은 봄베의 메뉴가 \texttt{A}--\texttt{B}--\texttt{C}--\texttt{A}의 루프 하나일 때, 어느
위치에서 생긴 연결을 그린 것이다. 루프를 펼쳐 그렸으므로 맨 오른쪽 열은 맨 왼쪽의 케이블 \texttt{A}와 같은 선이다.
그림을 단순하게 하려고 대각선 보드의 연결은 그리지 않았다.

\begin{figure}[ht]
\centering
\begin{tikzpicture}[x=1cm,y=-5.5mm,
    dot/.style={circle,draw,fill=white,inner sep=1.5pt}]
  \foreach \col/\name/\key in {0/A/A,3/B/B,6/C/C,9/A/Z} {
    \node[font=\small\bfseries] at (\col,-1.1) {\texttt{\name}};
    \foreach \l [count=\i from 0] in {a,b,c,d,e,f} {
      \node[dot] (\key\l) at (\col,\i) {};
    }
  }
  \foreach \l [count=\i from 0] in {a,b,c,d,e,f} {
    \node[font=\ttfamily\scriptsize,anchor=east] at (-0.15,\i) {\l};
  }
  \foreach \x/\y in {a/b,b/a,c/f,d/e,e/d,f/c} { \draw[wire] (A\x) -- (B\y); }
  \foreach \x/\y in {a/b,b/a,c/d,d/c,e/f,f/e} { \draw[wire] (B\x) -- (C\y); }
  \foreach \x/\y in {a/c,b/e,c/a,d/f,e/b,f/d} { \draw[wire] (C\x) -- (Z\y); }
  \node[note] at (1.5,6.2) {$\Scr_1$};
  \node[note] at (4.5,6.2) {$\Scr_2$};
  \node[note] at (7.5,6.2) {$\Scr_3$};
  \draw[cgray,dashed,->] (9,6.6) -- (9,7.1) -| node[pos=0.25,below,note] {같은 케이블 \texttt{A}} (0,6.6);
\end{tikzpicture}
\caption{작은 봄베의 루프 메뉴 \texttt{A}--\texttt{B}--\texttt{C}--\texttt{A}가 어느 위치에서 만드는 연결. 점 18개(케이블 셋 × 선
여섯)와 선분 18개로 된 그래프다. 다음 장에서 이 그림에 전압을 건다.}
\label{fig:c-graph}
\end{figure}

\Why 그래프로 보면 봄베의 둘째 일이 익숙한 물음이 된다. ``이 점에서 선분을 따라 갈 수 있는 점은 어디인가.''
그래프의 점들은 서로 이어진 \term{덩어리}{connected component}로 나뉘고, 전압을 건 선의 덩어리 전체가 그 가설의 귀결
전부다. 그림~\ref{fig:c-graph}의 18개 점도 덩어리 몇 개로 나뉜다. 어떻게 나뉘는지, 그리고 그것이 왜 해독의 답이
되는지가 다음 장의 내용이다.

\section{정리}

\begin{itemize}
\item 봄베는 위치 17,576가지마다 ``메뉴를 만족하는 플러그보드가 있을 수 있는가''를 묻는다.
\item 글자마다 짝에 대한 가설이 26개이고, 가설 하나가 선 하나다. 선은 676가닥이다.
\item 메뉴의 선 하나는 스크램블러를 지나 두 케이블의 선 26쌍을 잇는다. 연결 하나가 추론 하나이고, 스크램블러가
되돌릴 수 있으므로 방향이 없다.
\item 대각선 보드는 ``$X$의 짝이 $y$이면 $Y$의 짝은 $x$''를 325개의 고정 연결로 넣는다.
\item 메뉴의 연결은 위치마다 바뀐다. 봄베의 일은 \emph{연결을 만드는 것}과 \emph{이어진 선을 찾는 것} 둘이다.
\item 선과 연결은 그래프이고, 이어진 선을 찾는 것은 그래프의 덩어리를 찾는 것이다.
\end{itemize}
